
%regresa el valor de 12 pixeles
\normalsize
%\thispagestyle{empty}


\noindent \begin{center}
\textbf{Dedicatoria}
\par\end{center}

 \vspace{0.7mm}
 
 \begin{center}
     A mis padres, Carolina Duque y Andrés Giraldo, por ser el origen de todo lo que soy, aparte de su amor y el apoyo incondicional detrás de cada paso en este camino.
     
     A mi madrastra, Natalia Pérez, por su amor, su dedicación y su constante acompañamiento.
     
     A mi hermano, Matías Giraldo, por ser mi compañero de vida, por su apoyo y por los momentos que compartimos.
     
     A mis abuelos, Edelmira Murcia y Carlos Giraldo, Lucy Torrente y Alberto Duque, y a Patricia Marín, por sembrar en mí, con su ejemplo, el valor del esfuerzo y la constancia.
     
     A mis tías, Alejandra Duque, Martha Duque, Aura María Giraldo y Catalina Ríos, a mis tíos, Alberto Duque, Guillermo Gaviria, Julián Giraldo y Carlos Andrés Pérez, y a mi primo Jacobo Gaviria, por su cariño y por acompañarme siempre.
     
     A mis ahijados, Lucca y Benjamin Gaviria, que Dios les tenga preparada una vida llena de salud y prosperidad.
     
     A mi novia, Nycolle López, por su paciencia, su compañía y su aliento en los momentos más exigentes de este proceso.
     
     A Thomas Grisales, Simón Gaviria, Jacobo Ortiz y Daniel Pulgarin, mis amigos de toda la vida, por su amistad incondicional y por las historias que compartimos.
     
     A Jaime Miranda, Alejandro Vega y Juan José Mesa, amigos de estos últimos años, por su compañía sincera, por los momentos que hemos vivido juntos.
     
     A Jostin Gómez, por ser mi compañero inseparable en esta travesía académica, compartiendo cada aula, cada reto y cada logro que define nuestra formación universitaria.
     
     A todos ellos y a Dios, mi más profunda gratitud. \vspace{0.5cm}
 \end{center}
 
 \begin{flushright}
 	\textbf{Samuel Giraldo Duque}
 \end{flushright}
 
 \newpage
 
 \begin{center}
 	A quienes caminaron conmigo sin saber muy bien hacia dónde, y aun así no se detuvieron.
 	Hay una fuerza que no tiene nombre exacto y que, sin embargo, empuja a las personas a seguir, a preguntar, a no quedarse quietas ante lo desconocido. Esa fuerza —que corre por la sangre antes de tener explicación— fue la misma que me trajo hasta aquí. Ustedes la reconocieron en mí antes que yo, y la alimentaron con paciencia, incluso cuando el camino no tenía todavía forma ni nombre.
 	Esta tesis es, apenas, una de las respuestas posibles a esa pregunta que nunca dejamos de hacernos.
 	A Dios, que ha germinado en mí el don del entendimiento, dándome la fortaleza y la capacidad para seguir adelante.
 	A mi madre, Miryam Restrepo, por su amor incondicional y sostenerme incluso en los días más difíciles.
 	A mi padre, Henry Gómez, por enseñarme a caminar con paso firme y silencioso, y por estar ahí, siempre, aunque no siempre lo dijera con palabras.
 	A mi hermano, Cristian Gómez, mi primer compañero de vida, con quien comparto no solo la sangre sino la historia, las motivaciones, los sueños y las aspiraciones.
 	A mis abuelos maternos, Rocío Giraldo y Jesús Evelio Restrepo, por el cariño y el apoyo que se siente aún en la distancia y en el tiempo.
 	A mis abuelos paternos, María Vitalina Patiño y Miguel Ángel Gómez, quienes aunque la vida no me permitió disfrutarlos por mucho tiempo, dejaron una huella que indiscutiblemente también me formó.
 	A Juan David Quintana y Juan Sebastián Mosos, amigos que convirtieron el camino en algo más liviano de recorrer.
 	A Samuel Giraldo, compañero de tesis y, con el tiempo, algo más que un compañero: gracias por hacer de este proceso algo compartido, y no solo cumplido. \vspace{0.5cm}
 \end{center}
 
 \begin{flushright}
 	\textbf{Jostin Alejandro Gómez Restrepo}
 \end{flushright}


 %\vspace{6.0cm}
 \newpage

\noindent \begin{center} 
 \textbf{Agradecimientos}
\par\end{center}

 \vspace{0.7mm}
 
 \begin{center}
     Agradezco profundamente a mi familia por su apoyo incondicional, a mis amigos por su compañía y motivación constante, y a la Universidad de San Buenaventura por proporcionar los recursos necesarios para desarrollar este trabajo. Un agradecimiento especial al docente Jairo Iván Coy Coy, asesor de esta tesis, cuya guía experta y disponibilidad fueron fundamentales para la realización de este proyecto. \vspace{0.5cm}
 \end{center}
 
 \begin{flushright}
 	\textbf{Samuel Giraldo Duque} \vspace{1cm}
 \end{flushright}
 
 \begin{center}
 	Expreso mi sincero agradecimiento a la Universidad de San Buenaventura, por brindarme la formación académica, los recursos y el espacio necesarios para el desarrollo de este trabajo.
 	Un agradecimiento especial al docente Jairo Iván Coy Coy, asesor de esta tesis, cuya guía no se limitó a lo técnico: su disposición constante, su paciencia y su compromiso genuino con este proyecto marcaron una diferencia real en su desarrollo. Sus aportes y su exigencia fueron determinantes para llevar este trabajo a buen término.
 	Agradezco también a los docentes y funcionarios de la Universidad que, de una u otra forma, contribuyeron con su conocimiento y disposición al desarrollo de este trabajo.
 	A mi familia y amigos, gracias por su respaldo constante a lo largo de este proceso; su apoyo, aliento y comprensión fueron un soporte fundamental para poder cumplir con esta meta.
 	Finalmente, agradezco a todas las personas que, directa o indirectamente, aportaron su tiempo, conocimiento o colaboración para hacer posible este resultado. \vspace{0.5cm}
 \end{center}
 
 \begin{flushright}
 	\textbf{Jostin Alejandro Gómez Restrepo}
 \end{flushright}
 
 

\cleardoublepage{}
